% Ad Familiares 14.6 --- headnote
% ad-familiares-14-06, 14 July 48 BC, written at Dyrrhachium

\textit{Cicero to his household, written from Dyrrhachium on
the day before the Ides of July 48 BC (the manuscript
dateline: \textit{Scr. Dyrrhach i Id. Quint. a. 706 (48)};
the body, written at the moment of despatch, gives
``Idib. Quint.''). The salutation is the bare
\textit{Suis salutem dicit} --- ``to his own, greetings'' ---
not to Terentia by name, and the plural verbs in the body
(\textit{videatis}, \textit{scitis}) confirm a collective
address. The letter is the last surviving piece written
from Pompey's camp before Pharsalus (9 August); within a
month the campaign that has held Cicero in Epirus through
the spring will be lost on the plain in Thessaly.}

\textit{The content is purely practical and entirely
financial. A property has failed to sell; Cicero asks his
household to find some other way to satisfy a creditor
whose identity they already know. He acknowledges, with
characteristic dryness, that Tullia's thanks for some
unspecified kindness are well earned. And Pollex --- a
slave courier of his, whose tardiness keeps recurring in
these months --- is to be got on the road at once. The
register is hurried and elliptical: he has nothing to say,
no one to send letters by, and is making this one do the
work of a tally.}
